\input{./._relpath-to-latexroot.ltx}
\documentclass[\latexroot/\projectname]{subfiles}
\input{\latexroot/@resources/subfile-setup.ltx}
\begin{document}

\section{Conclusions}\whenintegrated{\label{conclusion}}
\whenintegrated{\label{sec:conclusion}} % integrated doc: SST for labels

The aftermath of the Global Financial Crisis reignited interest in the distributional consequences of aggregate shocks and the role that household heterogeneity plays in the amplification and propagation of said shocks. In this paper we constructed a Heterogeneous Agent New Keynesian Small Open Model Economy (HANKSOME) framework, specifically designed to study these challenges in the international macroeconomics context.

We applied our HANKSOME framework to study a current account expansion and subsequent reversal inspired by the experience of Hungary. We found that a fixed exchange rate regime would have significantly improved welfare in consumption equivalent terms, relative to a flexible exchange rate. This result obtains despite the fact that keeping the fixed exchange rate generates a recession. The desirability of a fixed versus floating exchange rate depends crucially on the distribution of household leverage, and the fraction of debt denominated in foreign currency.

% Smart bibliography: Only include bibliography if standalone AND has citations
\smartbib

\end{document}
